\title{Investigating Deep-UV Driven Photochemistry in Alkenes with Sub-10\,fs Time Resolution}
\author{Matt Bain}
\date{}

\begin{document}
\maketitle

\begin{abstract}
Deep-UV ultrafast spectroscopy has long been constrained by the pulse sources available to drive it: tunable sub-20\,fs UV light has typically meant complex (N)OPA setups, with pulse energy and duration locked in a trade-off. Resonant dispersive wave (RDW) generation in gas-filled hollow-core fibers offers a simpler alternative --- but until recently, the usable pulse energies were too low for many time-resolved experiments.

That has changed. Scaling the fiber core radius $a$ lets the loss length $L_\text{loss} \propto a^3$ outgrow the soliton fission length $L_\text{fiss} \propto a^2\tau$, preserving the self-compression dynamics that generate the RDW while opening the door to much higher pulse energies. The result is a single-fiber source that is continuously tunable from the near-IR to the VUV, delivers sub-10\,fs pulses, and now reaches microjoule-level energies --- a genuine step change in what deep-UV ultrafast sources can deliver.

I will describe this source and its application to gas-phase ultrafast X-ray scattering at LCLS, where sub-10\,fs deep-UV pulses are used to initiate and track the excited-state dynamics of 1,3-butadiene. The resulting data resolve the branching between competing conical-intersection pathways on the $S_1$ surface, revealing a scattering signature tied to the isomerization channel that time-averaged theoretical treatments had obscured --- and pointing to where current simulations underestimate the isomerization yield.
\end{abstract}

\end{document}
